\documentclass[../main.tex]{subfiles}

\begin{document}

% -----------------------------
\section{Installation}
% -----------------------------
If your Python version is >=3.10, you are good to go. Install fresh virtual environment as given:
\begin{lstlisting}[language=bash, caption={Setting up environment}]
mkdir xdsl-playground
cd xdsl-playground
python3 --version
python3 -m venv .venv
source .venv/bin/activate
\end{lstlisting}
After activation, your shell prompt should show (.venv) in front. Now, install xdsl:
\begin{lstlisting}[language=bash, caption={xDSL installation}]
pip install xdsl
xdsl-opt --help
echo 'builtin.module {}' | xdsl-opt
\end{lstlisting}
You should see a trivial IR printed back.
Let’s start with a very small hand-written IR file. Create \texttt{test.mlir}:
\begin{lstlisting}[language=Python, caption={xDSL example}]
builtin.module {
  func.func @add_two(%arg0 : i32, %arg1 : i32) -> i32 {
    %0 = arith.addi %arg0, %arg1 : i32
    return %0 : i32
  }
}
\end{lstlisting}
and run:
\begin{lstlisting}[language=bash, caption={}]
xdsl-opt test.mlir
\end{lstlisting}
If xDSL has that dialect set enabled, it should parse and print it.
The syntax is MLIR-like:
\begin{itemize}
    \item builtin.module is the top module.
    \item func.func is a function.
    \item arith.addi is an arithmetic op.
\end{itemize}
Each op has:
\begin{itemize}
    \item name (arith.addi)
    \item operands (\%arg0, \%arg1)
    \item result(s) (\%0)
    \item type annotation (: i32)
\end{itemize}

\end{document}
